% Drawing 142/06-14 -- folder 142, batch 06, archive page 112 (PDF page 113).
% Transcription: Un petit cercle avec les points $1$, $-1$ et $\zeta$, l'arc de $1$ à $\zeta$ fléché.
% Kind: circle with three points. Method 2 (schematic redraw from the facsimile). Status: ok.
% Notes: A circle with the points -1, 1 and zeta; the centre joined to 1 and to zeta; the arc from 1 to zeta (overinked) with an arrow at zeta.
% Coordinates: 1 unit = 1 mm = 1/100 of the crop width (bbox in page fractions: [0.18, 0.7, 0.38, 0.82]).
% Redrawn by Claude (claude-opus-5-5), issue #73, 2026-11-07; not checked by a human; not compiled (no TeX in the session).
% Rights: reproduces part of the work -- behind RENDER_TRANSCRIPTS, never published by itself.
\documentclass[tikz,border=3pt]{standalone}
\usepackage{amsmath,amssymb}
\usetikzlibrary{patterns,arrows.meta}
\tikzset{pencil/.style={black!45}, struck/.style={black!35}, mark/.style={circle,draw,inner sep=1.3pt}, pt/.style={circle,fill,inner sep=0.9pt}, opt/.style={circle,draw,fill=white,inner sep=1.1pt}}
\begin{document}
\begin{tikzpicture}[x=1mm,y=1mm,line cap=round,line join=round,font=\scriptsize]
\draw (51,-52) circle[radius=33];
\node[pt] at (19,-50) {};
\node[pt] at (83,-50) {};
\node[pt] at (36,-23) {};
\draw (50,-49) -- (83,-50);
\draw (50,-49) -- (36,-23);
\draw[-{Stealth[length=2mm]},very thick] plot[smooth, tension=0.55] coordinates {(83,-47) (80,-30) (65,-19) (48,-19) (37,-23)};
\node[inner sep=1pt] at (12,-51) {\scriptsize $-1$};
\node[inner sep=1pt] at (89,-52) {\scriptsize $1$};
\node[inner sep=1pt] at (30,-15) {\scriptsize $\zeta$};
\end{tikzpicture}
\end{document}
